% Retained preamble for this source, sliced MECHANICALLY from its own main TeX.
% Source file : source0.tex (lines 1-165 of the file)
% The block below is verbatim author text; the wrapper class, document metadata and
% \begin{document} are omitted because the assembler supplies them.
% lines taken: 1, 2, 3, 4, 5, 6, 7, 8, 9, 10, ... (165 lines total)
%%~Mouliné par MaN_auto v.0.32.0 2024-07-18 15:02:27
\documentclass[AHL,Unicode,longabstracts,published]{cedram}

\usepackage{enumerate}
\usepackage{tikz}
\usepackage{subfigure}




%\mathtoolsset{showonlyrefs}




\DeclareMathSymbol{\leqslant}{\mathalpha}{AMSa}{"36}
\DeclareMathSymbol{\geqslant}{\mathalpha}{AMSa}{"3E}
\DeclareMathSymbol{\eset}{\mathalpha}{AMSb}{"3F}
\renewcommand{\leq}{\;\leqslant\;}
\renewcommand{\geq}{\;\geqslant\;}
%\newcommand{\dd}{\,\text{\rm d}}




\newcommand{\la}{\label}
\newcommand{\Om}{\Omega}


\newcommand{\e}{\varepsilon}
\newcommand{\h}{\frac{1}{2}}

\newcommand{\rhoN}{\rho^{\pi}}
\newcommand{\si}{\sigma}
\newcommand{\tc}{\, |\, }

\newcommand{\scalar}[2]{\langle #1, #2\rangle}



\DeclareMathOperator*{\union}{\bigcup}
\DeclareMathOperator*{\inter}{\bigcap}
\DeclareMathOperator{\var}{Var}
\newcommand{\tv}{\mathrm{TV}}
\newcommand{\ent}{\mathrm{Ent}}
\newcommand{\gap}{\mathrm{gap}}
\newcommand{\rBe}{\mathrm{Be}}
\newcommand{\IND}{\mathbf{1}}
\newcommand{\ind}{\mathbf{1}}

\newcommand{\cA}{\mathcal{A}}
\newcommand{\cB}{\mathcal{B}}
\newcommand{\cE}{\mathcal{E}}
\newcommand{\cL}{\mathcal{L}}
\newcommand{\cP}{\mathcal{P}}
\newcommand{\cS}{\mathcal{S}}
\newcommand{\cT}{\mathcal{T}}
\newcommand{\cU}{\mathcal{U}}
\newcommand{\cV}{\mathcal{V}}
\newcommand{\bbC}{\mathbb{C}}
\newcommand{\bbE}{\mathbb{E}}

\newcommand{\bbN}{\mathbb{N}}
\newcommand{\bbP}{\mathbb{P}}
\newcommand{\bbR}{\mathbb{R}}
\newcommand{\bbS}{\mathbb{S}}


%\DeclareMathOperator{\cov}{Cov}
%\DeclareMathOperator{\Hess}{Hess}
%\DeclareMathOperator{\dist}{dist}
%\DeclareMathOperator{\osc}{osc}
%\DeclareMathOperator{\Id}{Id}
%\DeclareMathOperator{\diag}{diag}
%\DeclareMathOperator{\supp}{supp}
%\DeclareMathOperator{\PI}{PI}
%\DeclareMathOperator{\Ent}{Ent}




\let\a=\alpha
\let\b=\beta
\let\d=\delta
\let\e=\varepsilon
\let\g=\gamma
\let\l=\lambda
\let\o=\omega
\let\r=\rho
\let\t=\tau
\let\z=\zeta
\let\G=\Gamma
\let\O=\Omega





%%%%%%%%%%%%%%%%%%%%%%%%%%%%%%%%%%%%%%%%%%%%%%%%%%%%%%%%%%%%%%%%%%%%%%%%%%%%%%%%%%%%%%%%%%%%%%%%%%%
\graphicspath{{./figures/}}

\newcommand*{\mk}{\mkern -1mu}
\newcommand*{\Mk}{\mkern -2mu}
\newcommand*{\mK}{\mkern 1mu}
\newcommand*{\MK}{\mkern 2mu}

\hypersetup{urlcolor=purple, linkcolor=blue, citecolor=red}

\newcommand*{\romanenumi}{\renewcommand*{\theenumi}{\roman{enumi}}}
\newcommand*{\Romanenumi}{\renewcommand*{\theenumi}{\Roman{enumi}}}
\newcommand*{\alphenumi}{\renewcommand*{\theenumi}{\alph{enumi}}}
\newcommand*{\Alphenumi}{\renewcommand*{\theenumi}{\Alph{enumi}}}
\let\oldtilde\tilde
\renewcommand*{\tilde}[1]{\mathchoice{\widetilde{#1}}{\widetilde{#1}}{\oldtilde{#1}}{\oldtilde{#1}}}
\let\oldhat\hat
\renewcommand*{\hat}[1]{\mathchoice{\widehat{#1}}{\widehat{#1}}{\oldhat{#1}}{\oldhat{#1}}}
\let\oldforall\forall
\renewcommand*{\forall}{\mathrel{\oldforall}}
%%%%%%%%%%%%%%%%%%%%%%%%%%%%%%%%%%%%%%%%%%%%%%%%%%%%%%%%%%%%%%%%%%%%%%%%%%%%%%%%%%%%%%%%%%%%%%

\title[Nonlinear recombinations and generalized random transpositions]{Nonlinear recombinations and generalized random transpositions}

\alttitle{Recombinaisons non linéaires et transpositions aléatoires généralisées}


\subjclass{82C40, 82C20, 39B62}


\keywords{Nonlinear recombinations, Entropy, Kac program, Permutations, Logarithmic Sobolev inequalities}


\author[\initial{P.} \lastname{Caputo}]{\firstname{Pietro} \lastname{Caputo}}
\address{Department of Mathematics\\ and Physics,\\
Roma Tre University,\\
Largo San Murialdo 1,\\
00146 Roma (Italy)}
\email{pietro.caputo@uniroma3.it}


\author[\initial{D.} \lastname{Parisi}]{\firstname{Daniel} \lastname{Parisi}}
\address{Department of Mathematics\\
and Physics,\\
Roma Tre University,\\
Largo San Murialdo 1,\\
00146 Roma (Italy)}
\email{daniel.parisi@uniroma3.it}




\begin{abstract}
We study a nonlinear recombination model from population genetics as a combinatorial version of the Kac--Boltzmann equation from kinetic theory. Following Kac's approach, the nonlinear model is approximated by a mean field linear evolution with a large number of particles. In our setting, the latter takes the form of a generalized random transposition dynamics. Our main results establish a uniform in time propagation of chaos with quantitative bounds, and a tight entropy production estimate for the generalized random transpositions, which holds uniformly in the number of particles. As a byproduct of our analysis we obtain sharp estimates on the speed of convergence to stationarity for the nonlinear equation, both in terms of relative entropy and total variation norm.
\end{abstract}

\begin{altabstract}
Nous étudions un modèle de recombinaison non linéaire issu de la génétique des populations en tant que version combinatoire de l'équation de Kac--Boltzmann de la théorie cinétique. En suivant l'approche de Kac, le modèle non linéaire est approché par une évolution linéaire de champ moyen avec un grand nombre de particules. Dans notre cadre, ce dernier prend la forme d'une dynamique généralisée de transpositions aléatoires. Nos principaux résultats établissent une propagation du chaos uniforme dans le temps, avec des bornes quantitatives, et une estimation de la production d'entropie pour les transpositions aléatoires généralisées, qui est uniforme en le nombre de particules. En conséquence de notre analyse, nous obtenons des estimations précises sur la vitesse de convergence à l'équilibre pour l'équation non linéaire, tant en termes d'entropie relative que de norme de variation totale.
\end{altabstract}

\datereceived{2022-07-15}
\daterevised{2024-05-14}
\dateaccepted{2024-06-10}

\editors{S. Gou\"ezel and N. Privault}
%%%%%%%%%%%%%%%%%%%%%%%%%%%%%%%%%%%%%%%%%%%%%%%%%%%%%%%%%%%%%%%%%%%%%%%%%%%%%%%%%%%%%%%%%%%%%%%%%%%%%%%%%
\dateposted{2024-09-05}
